\section{Conclusion and Reproducibility}
\label{sec:conclusion}

\subsection{Conclusion}
This paper presented HOROS PROBE-1, a predictive Control Barrier Function safety framework with Adaptive Conformal Inference for safe multi-fingered robotic grasping under coupled nonstationary contact conditions. By coupling linearized Hunt-Crossley elastomeric dynamics with an 8-sided polyhedral friction cone inner approximation and grasp matrix wrench balance, the proposed controller provides real-time proactive force modulation within an embedded quadratic program. 

Across 3,400 simulation trials in Isaac Lab with PhysX 5 continuous collision physics, HOROS achieved a \textbf{0.0\% observed drop rate and 0.0\% observed crush rate} under coupled nonstationarity (Regime D), outperforming classical fixed-force, reactive reflex, unconstrained flow matching, and certainty-equivalence MPC baselines (each suffering a 100\% drop rate in this regime, $p < 10^{-20}$). Systematic 7-component ablations proved that slip prevention and structural object integrity are competing physical objectives that require two-sided barrier constraints coupled with full proprioceptive kinematics. The embedded C-level OSQP solver executes with a median reported latency of \textbf{0.193~ms}, and a maximum observed controller-cycle latency of \textbf{9.603~ms} in the tested configuration, remaining below the 10.0~ms (100~Hz) watchdog deadline during simulation runs. Future research will explore embedding this predictive safety filter as an inductive bias directly within diffusion-based foundation manipulation policies.

\subsection{Code and Data Availability}
To support independent verification and peer review, all source code, mathematical dynamics models, evaluation harnesses, and raw per-trial JSON experimental logs are open-sourced under an MIT License:
\begin{itemize}
    \item \textbf{Source Repository:} \url{https://github.com/Sunilkumarchavhan9/horos-probe-1}
    \item \textbf{Evaluation Engine:} \texttt{research/PROBE-1/experiments/run\_full\_evaluations.py}
    \item \textbf{Raw Data Ledger:} \texttt{research/PROBE-1/experiments/full\_evaluation\_results.json}
    \item \textbf{Unit Test Suite:} \texttt{research/PROBE-1/tests} ($26$ passing unit tests)
\end{itemize}

